\documentclass[10pt,a4paper]{amsart}
\usepackage[margin=19mm]{geometry}
\usepackage{amsmath,amssymb,mathtools}
\usepackage[scr=rsfs,cal=euler]{mathalfa}
\usepackage{xfrac}
\usepackage[unicode,hidelinks,bookmarks=false]{hyperref}
\newcommand{\ie}{i.e.\ }
\newcommand{\cf}{cf.\ }
\newcommand{\resp}{respectively\ }
\newcommand{\Subsection}{\subsection}
\providecommand{\nobreakdash}{\nobreak\hbox{-}}
\setlength{\emergencystretch}{2em}
\multlinegap0pt
\usepackage{amssymb}
\usepackage{mathtools}
\usepackage[shortlabels]{enumitem}
\usepackage{tikz}
\usepackage{dsfont}
\newtheorem{lemma}{Lemma}
\newtheorem{proposition}{Proposition}
\newtheorem{claim}{Claim}
\newtheorem{poc}{Poc}
\newcommand\dT{\mathcal{T}}
\renewcommand\ge{\geqslant}
\renewcommand\geq{\geqslant}
\renewcommand\le{\leqslant}
\renewcommand\leq{\leqslant}
\title{Triangle packings in randomly perturbed graphs}
\author{Xinbu Cheng, Hong Liu, Lanchao Wang, Zhifei Yan}
\date{Source: arXiv:2604.25250v2 (CC BY 4.0)}
\begin{document}
\maketitle

\noindent\emph{Source and licence.} Triangle packings in randomly perturbed graphs. Version arXiv:2604.25250v2 (CC BY 4.0), distributed by arXiv under the Creative Commons Attribution 4.0 International licence, http://creativecommons.org/licenses/by/4.0/, as recorded in the arXiv licence metadata for this exact version and shown on the abstract page https://arxiv.org/abs/2604.25250v2. Full licence notice: \texttt{licenses/arxiv-2604.25250-CC-BY-4.0.txt}.\par\smallskip

\noindent\textit{Editorial scope.} Counted unit: the article's lemma lem:regularity (source0.tex lines 359--363). The article's complete original proof of it is delivered with it (source0.tex lines 588--659, one \texttt{proof} environment of 72 lines, which the article places away from the statement). No mathematics is written, repaired or re-derived: the statement and the proof are the article's own lines, in the article's own notation, and no retained line is substituted or re-flowed.  Retained uncounted as the source-local context the proof needs: lines 368--384; lines 443--446; lines 465--480; lines 505--507; lines 509--512; lines 530--547.  Nothing was omitted from any retained span, so the package carries no omission record.  No retained line contains a citation, so the wrapper carries no bibliography and no bibliography entry.  No retained line contains a figure environment, an image-inclusion command or drawing code, so no image is delivered and the package declares no asset dependency.  Editorial formatting only, in addition: an AMS \texttt{amsart} preamble that restates the article's own declarations and macros used by the retained lines, namely the environments \texttt{lemma}, \texttt{proposition}, \texttt{claim}, \texttt{poc} on separate counters, together with the standard packages those lines need.  The retained environments print in this document as Lemma 1, Proposition 1, Claim 1, Poc 1 in order of appearance, whereas the article prints the same items with the numbering of its own full document.  The complete native source of the article as distributed in the arXiv e-print for this exact version is bundled unchanged as \texttt{sources/199290/source0.tex}.
\begin{lemma}[Triangle counting]\label{lem:trianglecounting}
For $\alpha,\beta,\gamma>p>0$, there exists $C=C(p)>0$ such that the following holds. For any $0<\varepsilon<p/2$, let $G[X,Y,Z]$ be a $3$-partite graph with $d(X,Y)=\alpha$, $d(Y,Z)=\beta$ and $d(Z,X)=\gamma$, and all the three pairs are $\varepsilon$-regular. Then there exists a collection of triangles $\mathcal T$ in $G$ such that
\[
|\mathcal T|
\ge
(1-C\varepsilon)\alpha\beta\gamma |X||Y||Z|.
\]
Moreover, let $\dT(e)\subset\dT$ be the set of triangles containing $e$ in $\dT$ for each edge $e\in E(G)$. Then
\begin{equation}\label{eq:Taue}
|\dT(e)| \le
\begin{cases}
(1+C\varepsilon)\beta\gamma |Z| & \text{if } e \in E(X,Y), \\[4pt]
(1+C\varepsilon)\alpha\gamma |X| & \text{if } e \in E(Y,Z), \\[4pt]
(1+C\varepsilon)\alpha\beta |Y| & \text{if } e \in E(X,Z).
\end{cases}
\end{equation}
\end{lemma}

\begin{proposition}\label{prop:fractionalpacking}
Let $d>0$ and $p> 2d/(1+2d)$, let $n$ be sufficiently large and let $G_d$ be a $dn$-regular graph on $n$ vertices. Then for any $\varepsilon>0$, with high probability there exists a fractional triangle packing $f$ of $G=G_d\cup G(n,p)$ such that
$$\|f\|\ge (1-\varepsilon)\cdot \frac{e(G)}{3}.$$
\end{proposition}

\begin{lemma}\label{lem:Gpartition}
Let $0<p,d\leq 1$ be constant, let $n$ be sufficiently large and let $G:=G_d\cup G(n,p)$ where $G_d$ is a $dn$-regular graph on $n$ vertices. Given $\varepsilon>0$, there exists $M\in\mathbb{N}$ such that the following holds. With high probability there exists a partition 
$$V(G)=V_1\cup...\cup V_M,\qquad\text{where}\ \big||V_i|-|V_j|\big|\leq 1\ \text{for all}\ i,j\in[M],$$ 
such that 
\begin{enumerate}
    \item[(i)] for all but at most $\varepsilon M^2$ pairs $(i,j)$ where $1\le i<j\le M$, the pair $(V_i,V_j)$ is $\varepsilon$-regular;

    \item[(ii)] for each $i\in [M]$, there are at most $\varepsilon M$ values of $j\in [M]$ such that $(V_i,V_j)$ is not $\varepsilon$-regular;

    \item[(iii)] for each pair $i,j\in [M]$, the corresponding edge density satisfies
    $$d(V_i,V_j)\geq (1-\varepsilon)p;$$

    \item[(iv)]  for each $i\in [M]$, the density sum 
    $$\sum_{j: j\neq i}d(V_i,V_j)\leq (1+\varepsilon)qM,$$ where $q: = p+(1-p)d$.
\end{enumerate}
\end{lemma}

\begin{align}\label{eq:di}
d_i\leq (1+\varepsilon)qM
\end{align}

\begin{align}\label{eq:W}
(1-4\varepsilon)q\binom{M}{2}\le W\le
(1+4\varepsilon)q\binom{M}{2}.   
\end{align}

\begin{lemma}\label{lem:reducedweight}
Let $\varepsilon>0$, let $M\in\mathbb{N}$ be sufficiently large and let $a,b,c>0$ satisfying $3a-3b/2+c/3>\varepsilon/M$. Let $R=K_M$ be a complete graph with $V(R)=\{v_1,...,v_M\}$, and equipped with an edge weight $w_{ij}$ for each $e_{ij}=v_iv_j\in E(R)$ satisfying the following properties:
\begin{enumerate}
    \item[$(W_1)$] for all but $\varepsilon M^2$ pair $i,j\in[M]$, the weight $w_{ij}\geq aM$;

    \item[$(W_2)$] for each $i\in[M]$, set $d_i:=\sum_{i\neq j}w_{ij}$, then $d_i\leq b\cdot\binom{M}{2}$;

    \item[$(W_3)$] the weight sum $ W:=\sum_{i<j} w_{ij}\geq c\cdot \binom{M}{3}$.
\end{enumerate}
Then there exists a triangle weight $x_{ijk}\geq 0$ for each triangle $T_{ijk}=\{e_{ij},e_{jk},e_{ik}\}\in\dT(R)$ where $1\leq i<j<k\leq M$ such that 
\begin{enumerate}
    \item [$(X_1)$] for each edge $e_{ij}$ $(1\leq i<j\leq M)$, the weight sum of triangles containing $e_{ij}$ is at most $w_{ij}$, that is
    $$\sum_{k\notin\{i,j\}}x_{ijk}\leq w_{ij};$$

    \item [$(X_2)$] the total sum of triangle weights satisfies
    $$\sum_{1\leq i<j<k\leq M}x_{ijk}\geq (1-4\varepsilon)\cdot \frac{c}{3}\cdot\binom{M}{3}.$$
\end{enumerate}
\end{lemma}

\begin{lemma}[Regularity]\label{lem:regularity}
For $0<\varepsilon<1$ and $m\ge 2$, there exists a constant $Z=Z(\varepsilon,m)$ such that the following holds. For every graph $G$, there exists $m\le M\le Z$ and a partition
$$V(G)=V_1\cup \cdots \cup V_{M} \qquad\text{where}\ \big||V_i|-|V_j|\big|\leq 1\ \text{for all}\ 1\leq i,j\leq M,$$
such that $(V_i,V_j)$ is $\varepsilon$-regular for all but at most $\varepsilon M^2$ pairs $(i,j)$ such that $1\le i<j\le M$. Moreover, for each $i\in [M]$, there are at most $\varepsilon M$ values of $j\in [M]$ such that $(V_i,V_j)$ is not $\varepsilon$-regular.
\end{lemma}

\begin{proof}[Proof of Proposition~\ref{prop:fractionalpacking}]
Given $G=G_d\cup G(n,p)$ with $d>0$ and $p>2d/(1+2d)$, set $q:=d+(1-d)p$. Let $\varepsilon>0$ be small enough. 

First, we apply regularity Lemma~\ref{lem:regularity} to $G$. By Lemma~\ref{lem:Gpartition}, there exists a partition $V(G)=V_1\cup...\cup V_M$ such that (i)-(iv) in Lemma~\ref{lem:Gpartition} hold, where $M\in\mathbb{N}$ is a sufficiently large constant. 

We next construct the reduced graph $R$ with respect to $G$. Let $R=K_M$ be a complete graph with $V(R)=\{v_1,...,v_M\}$. For each edge $e_{ij}:=v_iv_j$, set
$w_{ij}= d(V_i,V_j)$ if $(V_i,V_j)$ is an $\varepsilon$-regular pair in $G$. Otherwise set $w_{ij}=0$. Set $d_i:=\sum_{i\neq j}w_{ij}$ and $ W:=\sum_{i<j} w_{ij}$. By \eqref{eq:di} and \eqref{eq:W} we have
\begin{align}\label{eq:di'}
d_i\leq (1+\varepsilon)qM
\end{align}
holds for all $i\in [M]$, and
\begin{align}\label{eq:W'}
(1-4\varepsilon)q\binom{M}{2}\le W\le
(1+4\varepsilon)q\binom{M}{2}.   
\end{align}

Then we apply Lemma~\ref{lem:reducedweight} to $R$ to construct triangle weights for $R$. Set 
\begin{align*}
a=(1-\varepsilon)p/M,\qquad b=2(1+\varepsilon)q/M,\qquad c=3(1-4\varepsilon)q/M.   
\end{align*}
Note then ($W_1$) holds by (iii) of Lemma~\ref{lem:Gpartition}, and ($W_2$), ($W_3$) hold by \eqref{eq:di'} and \eqref{eq:W'} respectively. Moreover, by the choice of $a,b,c$ we have 
$$3a-3b/2+c/3\geq 3(1-\varepsilon)p/M - 3(1+\varepsilon)q/M + (1-4\varepsilon)q/M.$$
Since $p>2d/(1+2d)$ and $q=d+(1-d)p$, we have
$$3a-3b/2+c/3\geq \frac{1}{M}\cdot \Big(p-2d+2pd-\varepsilon\cdot(3p+7q)\Big)>\varepsilon/M$$
where the last inequality holds since $\varepsilon$ is small enough (we just need $\varepsilon<(p-2d+2pd)/(1+3p+7q)$). Therefore, by applying Lemma~\ref{lem:reducedweight} to $R$, there exists a triangle weight $x_{ijk}\geq 0$ for each triangle $T_{ijk}\in\dT(R)$ such that ($X_1$)-($X_2$) holds.  

Now let us construct triangle weights for each triangle in $G$. We first consider triangles between regular $3$-tuples. Let $1\leq i<j<k\leq M$ be an arbitrary 3-tuple such that $(V_i,V_j),(V_j,V_k)$ and $(V_i,V_k)$ are all $\varepsilon$-regular. Note then the corresponding density $w_{ij},w_{jk}, w_{ik}\geq(1-\varepsilon)p$ by (iii) of Lemma~\ref{lem:Gpartition}. Hence by choosing $\alpha=w_{ij}$, $\beta= w_{jk}$, $\gamma=w_{ik}$, we may apply Lemma~\ref{lem:trianglecounting} to the $3$-partite graph $G[V_i,V_j,V_k]$, to get a collection of triangles $\dT_{ijk}\subset \dT(G[V_i,V_j,V_k])$ such that
\begin{align}\label{eq:dTijk}
|\dT_{ijk}|\ge(1-C\varepsilon)\cdot\alpha\beta\gamma\cdot |V_i||V_j||V_k|\geq (1-C\varepsilon)\cdot w_{ij}w_{jk}w_{ik}\cdot(n/M)^3,   
\end{align}
and let $\dT_{ijk}(e)\subset\dT_{ijk}$ be the set of triangles containing $e$ in $\dT_{ijk}$ for $e\in E(G[V_i,V_j,V_k])$,
\begin{equation}\label{eq:Taue'}
|\dT_{ijk}(e)|\le
\begin{cases}
(1+C\varepsilon)\cdot w_{ik}w_{kj}\cdot (n/M) & \text{if } e \in E(V_i,V_j), \\[4pt]
(1+C\varepsilon)\cdot w_{ij}w_{jk}\cdot (n/M) & \text{if } e \in E(V_i,V_k), \\[4pt]
(1+C\varepsilon)\cdot w_{ij}w_{ik}\cdot(n/M) & \text{if } e \in E(V_j,V_k),
\end{cases}
\end{equation}
where $C>0$ is some absolute constant not relying on $\varepsilon$ and $i,j,k$, since $w_{ij},w_{jk}, w_{ik}\geq(1-\varepsilon)p>p/2$ for all regular 3-tuples $(i,j,k)$, and thus $C=C(p)$. For each triangle  $T\in\dT_{ijk}$, we set
\begin{align}\label{eq:fijk}
f(T)=f_{ijk}:=\frac{x_{ijk}}{(1+C\varepsilon)w_{ij}w_{jk}w_{ik}}\cdot\frac{M}{n}.   
\end{align}
For any other triangle $T\in\dT(G)$, we set $f(T)=0$. Then we claim such a weight function $f$ is a fractional triangle packing of $G$.

\begin{claim}
The function $f:\dT(G)\rightarrow [0,1]$ is a fractional triangle packing of $G$.
\end{claim}
\begin{poc}
For each edge $e\in E(G)$, it suffices to show the sum of the weights $f(T)$ over triangles $T$ containing $e$ is at most 1. Indeed, if $e$ is contained in some cluster $G[V_i]$ or $e\in (V_i,V_j)$ that is not an $\varepsilon$-regular pair, then $f(T)=0$ for all triangle $T$ containing $e$. Otherwise if $e\in E(V_i,V_j)$ that is $\varepsilon$-regular, the weight sum of triangles containing $e$ satisfies
$$\sum_{T\in \dT(G): e\in E(T)}f(T)\leq\sum_{k\notin\{i,j\}} f_{ijk}\cdot |\dT_{ijk}(e)|\leq(1+C\varepsilon)\cdot(n/M)\cdot\sum_{k\notin\{i,j\}}w_{ik}w_{kj}\cdot f_{ijk}$$
where the second inequality holds by \eqref{eq:Taue'} and $f(T)=0$ for all triangles in $\dT(G[V_i,V_j,V_k])$ but not in $\dT_{ijk}$. Further by \eqref{eq:fijk}, we may plug the value of $f_{ijk}$ to get
$$\sum_{T\in \dT(G): e\in E(T)}f(T)\leq \frac{1}{w_{ij}}\cdot\sum_{k\notin\{i,j\}}x_{ijk}\leq1$$
where the last inequality holds by ($X_1$) of Lemma~\ref{lem:reducedweight}, as required.
\end{poc}

Therefore, it remains to check $\|f\|\geq(1-\varepsilon)\cdot e(G)/3$. Indeed, by definition $\|f\|$ is the sum of all triangle weights in $G$, that is
\begin{align*}
\|f\|=\sum_{T\in\dT(G)}f(T)=\sum_{1\leq i<j<k\leq M}f_{ijk}\cdot|\dT_{ijk}|.
\end{align*}
By the estimation of $|\dT_{ijk}|$ (\eqref{eq:dTijk}) and the value of $f_{ijk}$ (\eqref{eq:fijk}), we have
\begin{align*}
\|f\|&=\sum_{T\in\dT(G)}f(T)=\sum_{1\leq i<j<k\leq M}f_{ijk}\cdot|\dT_{ijk}|\\
&\geq (1-C\varepsilon)\cdot \bigg(\frac{n}{M}\bigg)^3\cdot \bigg(\sum_{1\leq i<j<k\leq M}  w_{ij}w_{jk}w_{ik}\cdot \frac{x_{ijk}}{(1+C\varepsilon)w_{ij}w_{jk}w_{ik}}\cdot\frac{M}{n}\bigg)\\
&\geq \frac{(1-C\varepsilon)}{(1+C\varepsilon)}\cdot\bigg(\frac{n}{M}\bigg)^2\cdot\sum_{1\leq i<j<k\leq M}x_{ijk}\\
&\geq \frac{(1-C\varepsilon)}{(1+C\varepsilon)}\cdot\bigg(\frac{n}{M}\bigg)^2\cdot (1-4\varepsilon)\cdot \frac{c}{3}\cdot\binom{M}{3}\\
&\geq \frac{(1-C\varepsilon)}{(1+C\varepsilon)}\cdot\bigg(\frac{n}{M}\bigg)^2\cdot (1-4\varepsilon)\cdot \frac{3(1-4\varepsilon)q}{3M}\cdot\binom{M}{3}\\
&\geq (1-O(\varepsilon))\cdot \frac{q}{3}\cdot\binom{n}{2}\\
&\geq (1-O(\varepsilon))\cdot \frac{e(G)}{3},
\end{align*}
where the third inequality holds by ($X_2$) of Lemma~\ref{lem:reducedweight} and $c=3(1-4\varepsilon)q/M$, and the last inequality holds since $e(G)= (1\pm\varepsilon)q\binom{n}{2}$ holds with high probability, as required.
\end{proof}

\end{document}
